\documentclass{beamer}
\usepackage{SLaMLab_HHU_theme}

\usepackage{tikz}
\usepackage{grffile}

\usetikzlibrary{arrows,shapes}
\tikzstyle{every picture}+=[remember picture]
\everymath{\displaystyle}
\tikzstyle{na} = [baseline=-.5ex]

\usepackage[makeroom]{cancel}
\usepackage{booktabs}
%\usepackage{fontspec}
%\setmainfont{Doulos SIL}

\usepackage{microtype}

%\usepackage{xeCJK}
%\usepackage{CJKutf8}
%\begin{CJK*}{UTF8}{gbsn}

\usepackage{amssymb}
\usepackage{pifont}
\newcommand{\cmark}{\ding{51}}
\newcommand{\xmark}{\ding{55}}

%%%%%% Some other settings that can make things look nicer
%%% Set a smaller indent for description environment
\setbeamersize{description width=2em}
%%% Remove nav symbols (and shift any logo down to corner)
\setbeamertemplate{navigation symbols}{\vspace{-2ex}} 

\usepackage{appendixnumberbeamer}

\usepackage{comment}

\usepackage{multimedia}
\usepackage{media9}

%% Lines to help with generating handouts. Uncomment to activate
% \setbeamercolor{frametitle}{fg=black}
% \setbeamercolor{background canvas}{bg=white}
%%%%%%
% \usepackage{pgfpages}
% \pgfpagesuselayout{4 on 1}[letterpaper,landscape,border shrink=5mm]

\usepackage{hyphenat}

\title{TITLE OF TALK}
%\subtitle{Tackle the forgotten majority}
\author[shortname]{Kevin Tang\newline
	\newline\texttt{\footnotesize{kevin.tang@hhu.de}; \footnotesize{\url{slam.phil.hhu.de}}\\ \footnotesize{\url{www.kevintang.org}}}}
\date{\today}
\institute[shortinst]{Heinrich Heine University Düsseldorf}
\AtBeginSection{
\begin{frame}
\begin{center}
\structure{\Huge \insertsection}
\end{center}
\end{frame}
}

% Bibtex options
\def\printdateextralabel{\printlabeldateextra}

% Bib settings
\usepackage[backend=biber,
natbib=true,
%sorting=none,
bibstyle=biblatex-sp-unified,
citestyle=sp-authoryear-comp]{biblatex}

\newcommand{\cit}[1]{{\tiny #1}}

% Insert your bib file name here.
\bibliography{yourbibfile}

% Get rid of stupid beamer reference garbage like paper icon, and set text to black.
\setbeamertemplate{bibliography item}{}
\setbeamercolor{bibliography entry author}{fg=black}
\setbeamercolor{bibliography entry title}{fg=black}
\setbeamercolor{bibliography entry location}{fg=black}
\setbeamercolor{bibliography entry note}{fg=black}


% Custom Packages
\usepackage{subcaption}

\usepackage{eurosym}
\usepackage{dcolumn}
         
\usepackage[T1]{fontenc}
\usepackage{tipa}
\usepackage{vowel}
%	\usepackage{pst-vowel}
\usepackage{pifont}
\usepackage{graphicx}
\usepackage{textcomp}
%	\usepackage{tikz-qtree}
\usepackage{tabto}
\usepackage{setspace}
\usepackage{wasysym}
\usepackage{amsmath}
\usepackage{arydshln}
\usepackage{rotating}
\usepackage{multimedia}
%\usepackage{linguex}
\usepackage[normalem]{ulem}
\usepackage{xfrac}
\usepackage{booktabs}

% IPA-related commands
%
\newcommand{\ipa}[1]{\textipa{#1}}
\newcommand{\palj}{\super{j}}
\newcommand{\asp}{\ipa{\super{h}}}
\newcommand{\ej}{\ipa{\super{P}}}
\newcommand{\unre}{\textcorner}
\newcommand{\er}{@\textrhoticity}
\newcommand{\Er}{3\textrhoticity}
\newcommand{\darkl}{\textsuperimposetilde{l}}
\newcommand{\beltl}{\textbeltl}
% Also \*l
\newcommand{\paln}{\textltailn}
% Also \*n

% Phono-related example formatting commands
%
\newcommand{\Rank}[2]{{\sc #1} $\gg$ {\sc #2}}
\newcommand{\RankHG}[2]{\emph{w}({\sc #1}) $>$ \emph{w}({\sc #2})}
\newcommand{\w}[1]{\emph{w}({\sc #1})}
%\newcommand{\pdata}[2]{\textipa{[ #1 ]} & {`#2'}}
\newcommand{\UR}[1]{\textipa{/#1/}}
\newcommand{\SR}[1]{\textipa{[#1]}}
\newcommand{\IR}[1]{\textipa{$\vert$ #1 $\vert$}}
\newcommand{\spe}[3]{\ipa{#1} \rar \ipa{#2} / \ipa{#3} }
\newcommand{\rar}{$\rightarrow$~}
%\newcommand{\fder}[2]{\Fmat{#1} \rar \Fmat{#2}}
\newcommand{\F}[1]{{\sc [#1]}}
\newcommand{\foc}{\underline{~~~}}
\newcommand{\Fmat}[1]{$\left[\begin{array}{c}#1\end{array}\right]$}
\newcommand{\angbr}[1]{$\langle$#1$\rangle$}
\newcommand{\disjunct}[1]{$\left\{\begin{array}{c}#1\end{array}\right\}$}

\newcommand{\nrm}[1]{\textnormal{#1}}

\newcommand{\skipline}{\vspace{\baselineskip}}
\newcommand{\newL}{\skipline\noindent}
\newcommand{\flag}[1]{{\red{\textbf{#1}}}}

\newcommand{\citepos}[1]{\citeauthor{#1}'s (\citeyear{#1})}


\newcommand{\smc}[1]{{\rmfamily{\textsc{#1}}}}

\newcommand{\specialcell}[2][c]{%
  \begin{tabular}[#1]{@{}c@{}}#2\end{tabular}}
  
\addtobeamertemplate{frametitle}{\vskip0ex}{}
%\setbeamertemplate{frametitle}{%
%    \nointerlineskip%
%    \begin{beamercolorbox}[wd=\paperwidth]{frametitle}
%        \insertframetitle%
%    \end{beamercolorbox}%
%}


\begin{document}
\definecolor{Grey}{rgb}{0.8, 0.25, 0.33}
\definecolor{brickred}{rgb}{0.8, 0.25, 0.33}
\definecolor{burgundy}{rgb}{0.5, 0.0, 0.13}
\definecolor{carnelian}{rgb}{0.7, 0.11, 0.11}
\definecolor{carmine}{rgb}{0.59, 0.0, 0.09}
\newcommand{\newemph}[1]{{\color{carmine}{{#1}}}}
\frame{
  \titlepage
}




% \begin{frame}
% \tableofcontents[ 
% % currentsubsection, 
% hideallsubsections, 
% sectionstyle=show
% %, 
% %subsectionstyle=shaded, 
% ] 
% \end{frame}

% A flagging system for comments
\newcommand{\flagKT}[1]{{\color{blue}{\textbf{#1}}}}

\frame{
\frametitle{Outline}
\tableofcontents[ 
%currentsubsection, 
hideallsubsections, 
sectionstyle=show
%, 
%subsectionstyle=shaded, 
] 
}

\section{Demo section}

\begin{frame}{Demo: Slide title with citations}
\begin{itemize}
\item Normal citation \citep{TangShaw_2021_ProsodyCognition}
\item Mini citation \cit{\citep{TangShaw_2021_ProsodyCognition}}
\end{itemize}
\end{frame}

\begin{frame}{Demo: Slide title with one figure}
\includegraphics[width=0.7\linewidth]{images/WAG.png}\\
\end{frame}


\begin{frame}{Demo: Slide with two figures}

\begin{figure}
\centering
\begin{subfigure}{.5\textwidth}
  \includegraphics[width=0.7\linewidth]{images/WAG.png}
  \caption{Caption for figure 1}
\end{subfigure}%
\centering
\begin{subfigure}{.5\textwidth}
  \includegraphics[width=0.7\linewidth]{images/WAG.png}
  \caption{Caption for figure 2}
\end{subfigure}
\end{figure}

\end{frame}

\begin{frame}{Demo: Slide title with special emphasis}
\begin{itemize}
\item Emphasis \newemph{Hello}
\end{itemize}
\end{frame}

\section{Introduction}
\begin{frame}{Introduction}
\end{frame}

\section{Method}
\begin{frame}{Method}
\end{frame}

\section{Results}
\begin{frame}{Results}
\end{frame}

\section{Discussion}
\begin{frame}{Discussion}
\end{frame}

\section{Conclusions}
\begin{frame}{Conclusions}
\end{frame}



\end{document}
